% ============================================================================
%  QuestFlow — Flutter Coursework Documentation Report
%
%  HOW TO USE THIS FILE
%  1. Upload report.tex to https://www.overleaf.com (New Project → Blank,
%     or upload this file). Compile with pdfLaTeX (default).
%  2. Take the screenshots listed in Section 7 and upload them to a folder
%     named "figures" in the same Overleaf project. The \includegraphics
%     lines are already written — just uncomment them once the images are
%     uploaded (or keep them commented and add images one by one).
%  3. Fill in your name, roll number and section on the title page.
% ============================================================================

\documentclass[12pt,a4paper]{article}

\usepackage[margin=2.4cm]{geometry}
\usepackage{graphicx}
\usepackage{xcolor}
\usepackage{booktabs}
\usepackage{tabularx}
\usepackage{float}
\usepackage{enumitem}
\usepackage{listings}
\usepackage{caption}
\usepackage[hidelinks]{hyperref}

% ---- colours used by code listings ----------------------------------------
\definecolor{codebg}{RGB}{248,247,242}
\definecolor{codekw}{RGB}{90,42,42}
\definecolor{codestr}{RGB}{47,105,60}
\definecolor{codecomment}{RGB}{110,110,100}
\definecolor{rulecol}{RGB}{200,200,190}

% ---- Dart language definition for listings ---------------------------------
\lstdefinelanguage{Dart}{
  keywords=[1]{abstract,as,assert,async,await,break,case,catch,class,const,
    continue,default,do,else,enum,extends,extension,factory,false,final,
    finally,for,get,if,implements,import,in,is,late,library,mixin,new,null,
    on,operator,part,required,rethrow,return,set,static,super,switch,this,
    throw,true,try,typedef,var,void,while,with,yield},
  keywords=[2]{int,double,bool,String,List,Map,DateTime,Duration,Future,void},
  sensitive=true,
  comment=[l]{//},
  morestring=[b]',
  morestring=[b]",
}
\lstdefinelanguage{YAML}{
  keywords=[1]{name,on,push,branches,jobs,steps,uses,with,run,needs,permissions},
  sensitive=true,
  comment=[l]{\#},
  morestring=[b]',
  morestring=[b]",
}
\lstdefinelanguage{SQL}{
  keywords=[1]{create,table,alter,policy,enable,row,level,security,select,
    insert,update,delete,using,check,not,null,default,references,primary,key,
    for,to,authenticated,on,if,exists,constraint,trigger,before,each,row,
    function,execute,uuid,text,timestamptz},
  sensitive=false,
  comment=[l]{--},
  morestring=[b]',
}

\lstset{
  basicstyle=\ttfamily\footnotesize,
  backgroundcolor=\color{codebg},
  keywordstyle=\color{codekw}\bfseries,
  keywordstyle=[2]\color{blue!60!black},
  commentstyle=\color{codecomment}\itshape,
  stringstyle=\color{codestr},
  frame=single,
  rulecolor=\color{rulecol},
  breaklines=true,
  showstringspaces=false,
  tabsize=2,
  captionpos=b,
  xleftmargin=1em,
  framexleftmargin=0.6em,
}

\captionsetup{font=small,labelfont=bf}
\setlist{itemsep=2pt,topsep=4pt}

\begin{document}

% ============================================================================
% TITLE PAGE
% ============================================================================
\begin{titlepage}
  \centering
  \vspace*{2cm}
  {\Large \textbf{[University / College Name]}\\[4pt]
  \large [Department Name]\\[1cm]}
  \rule{\linewidth}{0.6pt}\\[0.4cm]
  {\huge \textbf{QuestFlow}\\[0.4cm]
  \Large A Responsive Flutter Web Application for\\ Tasks, Events and Continuous Quests\\[0.2cm]}
  \rule{\linewidth}{0.6pt}\\[1.2cm]

  {\large \textbf{Flutter Application Development — Single Credit Coursework}\\[2pt]
  Documentation Report\\[1.5cm]}

  % TODO: fill in your details
  \begin{tabular}{rl}
    \textbf{Submitted by} & [Your Name] \\[6pt]
    \textbf{Roll Number} & [Your Roll Number] \\[6pt]
    \textbf{Section}     & [Your Section] \\[6pt]
    \textbf{Date}        & November 2026 \\
  \end{tabular}

  \vfill
  {\large Live deployment: \url{https://sahithi-101.github.io/questflow/}}
  \vspace{1cm}
\end{titlepage}

\tableofcontents
\newpage

% ============================================================================
\section{Application Name and Introduction}
% ============================================================================

\textbf{QuestFlow} is a fully functional, responsive web application built
with Flutter for managing three kinds of personal activity:

\begin{itemize}
  \item \textbf{Tasks} — work that must be completed before a deadline
        (pending, completed, or automatically derived overdue states).
  \item \textbf{Events} — occurrences at a specific time, with optional end
        time, location and reminder offsets.
  \item \textbf{Quests} — continuous daily habits or measurable goals
        tracked on a heatmap, with streaks and completion statistics.
\end{itemize}

The application presents exactly three primary destinations --- Home,
Tasks \& Events, and Quests --- plus a fourth page, a \textbf{Focus Timer},
which records named time logs of work sessions. Every item supports full
create, read, update and delete (CRUD) operations. User accounts are managed
by Supabase authentication, and all records are stored in a PostgreSQL
database protected by Row Level Security, so each account sees only its own
data. The compiled application is deployed on GitHub Pages and can be used
from any browser, on desktop or mobile.

\begin{figure}[H]
  \centering
  % \includegraphics[width=0.85\textwidth]{figures/home-desktop.png}
  \caption{QuestFlow home dashboard (desktop width).}
  \label{fig:home}
\end{figure}

% ============================================================================
\section{Purpose and Objectives}
% ============================================================================

\subsection{Purpose}
QuestFlow was designed as a personal productivity application that combines
a compact daily dashboard, a calendar for time-bound items, and a
heatmap-based habit tracker. The purpose of the project is twofold:

\begin{enumerate}
  \item To build a genuinely useful day-to-day productivity tool, and
  \item To demonstrate, in one coherent product, the Flutter concepts
        covered during the laboratory sessions --- widgets, responsive
        design, navigation, animations, state management, theming, custom
        widgets, and forms with validation.
\end{enumerate}

\subsection{Objectives}
\begin{itemize}
  \item Provide exactly three primary destinations (Home, Tasks \& Events,
        Quests) with clear navigation, plus a focus timer.
  \item Implement complete CRUD operations for tasks, events, quests and
        time logs, with validation on every form.
  \item Render quest history as a heatmap implementing documented colour
        rules (grey / green gradient / full green / purple) and compute
        streaks and completion rates from it.
  \item Persist data in the cloud with per-user isolation enforced by
        Row Level Security in the database itself.
  \item Keep every screen responsive from mobile widths ($\approx$390\,px)
        to desktop widths ($\geq$1280\,px).
  \item Verify correctness with an automated test suite (81 tests).
  \item Deploy the release build continuously to GitHub Pages.
\end{itemize}

% ============================================================================
\section{Tools and Technologies Used}
% ============================================================================

\begin{table}[H]
\centering
\begin{tabularx}{\textwidth}{lX}
\toprule
\textbf{Tool / Technology} & \textbf{Role in the project} \\
\midrule
Flutter 3.47 (stable) & UI toolkit; one codebase compiled for web, Android and Linux desktop \\
Dart 3 & Programming language for all application and test code \\
go\_router (\texttt{\^{}14}) & Declarative navigation: shell route, nested routes, query parameters \\
supabase\_flutter (\texttt{\^{}2.8}) & Cloud backend client: authentication and PostgreSQL CRUD \\
Supabase Auth & Email/password sign-up, sign-in, password reset, persistent sessions \\
PostgreSQL + Row Level Security & Cloud storage; per-user data isolation enforced in the database \\
flutter\_test & Automated unit and widget testing (81 tests) \\
flutter\_launcher\_icons & Generation of launcher icons for Android, iOS, web and macOS \\
Git + GitHub & Version control and source hosting \\
GitHub Actions + GitHub Pages & Continuous deployment: tests, release build and hosting \\
VS Code / Android Studio & Development environment \\
Chrome and Android device & Testing on different screen sizes and platforms \\
\bottomrule
\end{tabularx}
\caption{Tools and technologies.}
\end{table}

\subsection{Architecture}
The application follows a layered architecture so that the UI never talks
to the storage backend directly:

\begin{center}
\begin{tabular}{c l l}
\toprule
Layer & Responsibility & Key files \\
\midrule
UI & Screens and reusable widgets & \texttt{lib/screens/}, \texttt{lib/widgets/} \\
State & Single \texttt{ChangeNotifier} store & \texttt{lib/state/app\_state.dart} \\
Data & Repository interfaces + 2 implementations & \texttt{lib/data/} \\
Cloud & PostgreSQL, Auth, RLS & Supabase project \\
\bottomrule
\end{tabular}
\end{center}

The in-memory implementation backs demo mode; the Supabase implementation
backs cloud mode. Both satisfy the same interfaces, so the whole storage
backend is swappable without touching a single screen.

% ============================================================================
\section{Features and Functionalities Implemented}
% ============================================================================

\subsection{Home dashboard}
\begin{itemize}
  \item Tasks summary bar: pending count, due today, overdue count.
  \item Events summary bar: events of today and tomorrow, chronological.
  \item Quests bar: completed today / active total with completion percentage.
  \item Every bar navigates to its screen; quick-access button to add items.
\end{itemize}

\subsection{Tasks \& Events (calendar page)}
\begin{itemize}
  \item Month calendar (Monday-first) with shape-and-colour markers:
        squares for tasks (filled = pending, outlined = done) and circles
        for events, with a legend.
  \item Month navigation, ``Today'' button, day selection.
  \item Agenda for the selected day, combining tasks, events and time logs
        in one chronological list.
  \item Separate task and event forms with validation; a single ``+'' flow
        asks which one to create.
  \item Detail dialogs for both kinds with edit and delete-with-confirmation.
  \item Tasks can be marked done and reopened; overdue is derived
        automatically.
\end{itemize}

\subsection{Quests}
\begin{itemize}
  \item Quest list with per-quest progress, current streak and completion rate.
  \item Quest form: binary or numeric tracking mode, unit, daily target,
        upper limit, start/end dates with month/year presets; the end-date
        picker cannot choose a date before the start date.
  \item Full heatmap screen: month blocks, colour rules --- grey for zero,
        green gradient for partial progress, full green at target, purple
        for over-target days; future days are visually distinct and never
        shown as failures.
  \item Day dialog to record or edit any past day; future days are blocked;
        values above the upper limit are rejected with a message, never
        silently capped.
  \item Statistics: current streak, longest streak, completion rate,
        successful days, elapsed days, overload days.
  \item Archive as a safer alternative to deletion.
\end{itemize}

\subsection{Focus Timer (fourth page)}
\begin{itemize}
  \item Stopwatch with Start, Pause/Resume, and End \& log; the session
        keeps running while the user visits other pages.
  \item On ending, the app asks for a name; empty names cannot be saved and
        the session can be discarded.
  \item Paused time is excluded from the log: ending while paused closes
        the log at the pause moment.
  \item Today's logs are listed with name, start--end time, duration and a
        daily total; each log also appears on the calendar agenda under its
        day and can be deleted from there.
  \item Logs are retained for 30 days from their calendar day.
\end{itemize}

\subsection{Authentication and cloud storage}
\begin{itemize}
  \item Email/password sign-up, sign-in, sign-out and password reset.
  \item Session restoration: returning users stay signed in.
  \item Row Level Security policies on \emph{every} table guarantee that
        one account cannot read or modify another account's records.
  \item Loading, error, empty and success states for all cloud operations.
\end{itemize}

\subsection{Data retention}
Finished items are cleaned up automatically by a sweep that runs on every
load and hourly: completed tasks one week after completion, events one week
after they end, quests one week after their end date (with their daily
records), and time logs 30 days after their calendar day.

\subsection{Coursework concept mapping}
\begin{table}[H]
\centering
\small
\begin{tabularx}{\textwidth}{lX}
\toprule
\textbf{Required concept} & \textbf{Where it is implemented} \\
\midrule
Different widgets & \texttt{Container}, \texttt{Row}, \texttt{Column},
\texttt{Stack} (live badge over the stopwatch face), \texttt{ListView},
\texttt{GridView.count} (calendar and heatmap grids), styled cards,
\texttt{FilledButton}/\texttt{OutlinedButton}/icon buttons,
\texttt{AlertDialog}, \texttt{TextField}, chips, \texttt{FloatingActionButton} \\
\midrule
Responsive UI & \texttt{MediaQuery.sizeOf} drives the sidebar$\leftrightarrow$bottom-bar switch at 860\,px and per-screen paddings; \texttt{LayoutBuilder} switches the calendar between one and two columns at 760\,px; \texttt{MediaQuery.disableAnimationsOf} respects reduced-motion settings \\
\midrule
Navigation & \texttt{go\_router} with a \texttt{ShellRoute} (shared shell), five top-level routes, a nested quest-detail route, and query-parameter-driven flows \\
\midrule
Animations & Fade+slide page transitions (\texttt{CustomTransitionPage}), \texttt{AnimatedContainer} on the stopwatch face and calendar day selection, \texttt{AnimatedSwitcher} on the timer status, plus implicit dialog and snackbar transitions \\
\midrule
State management & \texttt{StatefulWidget} throughout the UI; a single \texttt{ChangeNotifier} (\texttt{AppState}) exposed through a custom \texttt{InheritedWidget} (\texttt{AppScope}) --- the same provider pattern the \texttt{provider} package formalises \\
\midrule
Themes \& custom widgets & Custom ``Lichen \& Madder'' palette and \texttt{ThemeData}; reusable custom widgets: summary bars, agenda rows, status views, form dialog scaffold, confirm-delete dialog, heatmap cells \\
\midrule
Forms and validation & Task, event and quest forms with required-field checks, date rules (end date cannot precede start date), numeric range validation against the quest's upper limit, and non-empty log names \\
\bottomrule
\end{tabularx}
\caption{Mapping of the required coursework concepts to the implementation.}
\label{tab:concepts}
\end{table}

% ============================================================================
\section{Important Code Snippets}
% ============================================================================

\subsection{Central state management --- \texttt{ChangeNotifier}}
\texttt{AppState} holds every list the UI shows. Mutations go through it and
end with \texttt{notifyListeners()}, which rebuilds every listening widget.
\lstinputlisting[language=Dart, caption={AppState: the single store and its
notification pattern (lib/state/app\_state.dart)}]{snippets/app_state_mutation.dart}

\subsection{Repository pattern --- one interface, two backends}
Screens depend only on the interface; demo mode uses in-memory lists and
cloud mode uses Supabase.
\lstinputlisting[language=Dart, caption={Repository interface
(lib/data/repositories.dart)}]{snippets/repository_interface.dart}
\lstinputlisting[language=Dart, caption={Supabase implementation: create
returns the stored row so the app uses the database-generated id
(lib/data/supabase\_repositories.dart)}]{snippets/supabase_create.dart}

\subsection{Responsive layout}
\lstinputlisting[language=Dart, caption={Responsive shell: sidebar on wide
screens, bottom navigation bar on narrow ones
(lib/widgets/app\_shell.dart)}]{snippets/responsive_shell.dart}
\lstinputlisting[language=Dart, caption={Responsive calendar: one or two
columns depending on measured width (lib/screens/calendar\_screen.dart)}]{snippets/responsive_calendar.dart}

\subsection{Navigation with go\_router}
\lstinputlisting[language=Dart, caption={Route table with shell route and
nested quest detail route (lib/router.dart)}]{snippets/router.dart}

\subsection{Animations}
\lstinputlisting[language=Dart, caption={Fade-and-slide page transition
(lib/router.dart)}]{snippets/animation_page.dart}
\lstinputlisting[language=Dart, caption={AnimatedContainer recolours the
stopwatch face by session state; AnimatedSwitcher cross-fades the status
label (lib/screens/timer\_screen.dart)}]{snippets/animation_timer.dart}

\subsection{Forms and validation}
\lstinputlisting[language=Dart, caption={Task form validation: required
title and deadline (lib/widgets/forms/task\_form.dart)}]{snippets/form_validation.dart}
\lstinputlisting[language=Dart, caption={Quest day validation: future days
blocked, upper limit enforced, nothing silently capped
(lib/state/app\_state.dart)}]{snippets/quest_validation.dart}

\subsection{Custom theme and reusable widgets}
\lstinputlisting[language=Dart, caption={The ``Lichen \& Madder'' palette:
semantic item colours and the heatmap colour function
(lib/theme/palette.dart)}]{snippets/palette.dart}

\subsection{Database security --- Row Level Security}
\lstinputlisting[language=SQL, caption={Row Level Security on the time\_logs
table: every table in the project carries the same four policies
(supabase/migration\_002\_time\_logs.sql)}]{snippets/rls.sql}

\subsection{Testing}
\lstinputlisting[language=Dart, caption={One of 81 automated tests: the
pause/resume accumulation rules of the focus timer
(test/phase5\_test.dart)}]{snippets/test_example.dart}

% ============================================================================
\section{Deployment Details}
% ============================================================================

\subsection{Build and hosting}
The application is deployed to \textbf{GitHub Pages} at
\begin{center}
\url{https://sahithi-101.github.io/questflow/}
\end{center}

Deployment is fully automatic: every push to the \texttt{main} branch
triggers a GitHub Actions workflow that (1) installs Flutter, (2) runs the
complete test suite, (3) builds the release web bundle with the correct
base path for the repository URL, and (4) publishes \texttt{build/web} to
GitHub Pages. The workflow can also be triggered manually.

\lstinputlisting[language=YAML, caption={.github/workflows/deploy.yml}]{snippets/deploy.yml}

\subsection{Cloud backend}
\begin{itemize}
  \item A Supabase project hosts the PostgreSQL database and email/password
        authentication.
  \item The schema (\texttt{supabase/schema.sql} plus two migration files)
        defines six tables --- profiles, tasks, events, quests,
        quest\_daily\_records and time\_logs --- with RLS policies,
        \texttt{updated\_at} triggers, a unique constraint on
        \texttt{(quest\_id, record\_date)} and check constraints.
  \item Only the public anon key is embedded in the client; authorisation
        is enforced by RLS in the database itself.
  \item The deployed URL is registered in Supabase's allowed redirect URLs
        so confirmation and password-reset emails work from the live site.
\end{itemize}

\subsection{Running locally}
\begin{lstlisting}[language=, basicstyle=\ttfamily\small]
flutter pub get
flutter run -d chrome          # web (demo mode without Supabase keys)
flutter run -d linux           # Linux desktop build
flutter run -d android         # Android device over USB
flutter test                   # 81 automated tests
flutter build web --release    # production bundle
flutter build apk --release    # Android installer
\end{lstlisting}

% ============================================================================
\section{Screenshots of the Application}
% ============================================================================

% Every figure below expects one PNG in the figures/ folder. Suggested
% capture sizes: desktop = 1280x800 browser window; mobile = 390x844
% (Chrome DevTools device toolbar). Uncomment the \includegraphics lines
% as you add the images.

\subsection{Home dashboard}

\begin{figure}[H]
  \centering
  % \includegraphics[width=0.9\textwidth]{figures/home-desktop.png}
  \caption{Home dashboard on desktop: the three summary bars with live
  counts, quick-add button and sidebar navigation.}
\end{figure}

\begin{figure}[H]
  \centering
  % \includegraphics[width=0.32\textwidth]{figures/home-mobile.png}
  \caption{Home dashboard at mobile width: the sidebar is replaced by a
  bottom navigation bar and the bars stack vertically.}
\end{figure}

\subsection{Tasks \& Events}

\begin{figure}[H]
  \centering
  % \includegraphics[width=0.9\textwidth]{figures/calendar-desktop.png}
  \caption{Calendar with task/event markers and the selected-day agenda
  (two-column desktop layout).}
\end{figure}

\begin{figure}[H]
  \centering
  % \includegraphics[width=0.32\textwidth]{figures/calendar-mobile.png}
  \caption{Calendar at mobile width: single column, bottom navigation.}
\end{figure}

\begin{figure}[H]
  \centering
  % \includegraphics[width=0.55\textwidth]{figures/task-form-validation.png}
  \caption{Task form with validation errors shown after submitting empty
  required fields.}
\end{figure}

\begin{figure}[H]
  \centering
  % \includegraphics[width=0.55\textwidth]{figures/event-form.png}
  \caption{Event form: start/end pickers, location and reminder offsets.}
\end{figure}

\begin{figure}[H]
  \centering
  % \includegraphics[width=0.55\textwidth]{figures/task-detail.png}
  \caption{Task detail dialog with Edit and Delete actions.}
\end{figure}

\begin{figure}[H]
  \centering
  % \includegraphics[width=0.45\textwidth]{figures/delete-confirm.png}
  \caption{Delete confirmation dialog.}
\end{figure}

\begin{figure}[H]
  \centering
  % \includegraphics[width=0.9\textwidth]{figures/agenda-with-log.png}
  \caption{Agenda showing a focus-timer log (amber) among tasks and events,
  with its name, start--end time and duration.}
\end{figure}

\subsection{Quests}

\begin{figure}[H]
  \centering
  % \includegraphics[width=0.9\textwidth]{figures/quest-list.png}
  \caption{Quest list with per-quest streaks and completion rates.}
\end{figure}

\begin{figure}[H]
  \centering
  % \includegraphics[width=0.55\textwidth]{figures/quest-form.png}
  \caption{Quest form: tracking mode, unit, target, upper limit and date
  range with presets.}
\end{figure}

\begin{figure}[H]
  \centering
  % \includegraphics[width=0.9\textwidth]{figures/quest-heatmap.png}
  \caption{Quest heatmap: grey/partial-green/full-green/purple colour rules
  with future days visually distinct.}
\end{figure}

\begin{figure}[H]
  \centering
  % \includegraphics[width=0.45\textwidth]{figures/quest-day-dialog.png}
  \caption{Day dialog recording progress for one quest day.}
\end{figure}

\begin{figure}[H]
  \centering
  % \includegraphics[width=0.55\textwidth]{figures/quest-stats.png}
  \caption{Quest statistics: current streak, longest streak, completion
  rate, overload days.}
\end{figure}

\subsection{Focus Timer}

\begin{figure}[H]
  \centering
  % \includegraphics[width=0.5\textwidth]{figures/timer-idle.png}
  \caption{Timer page in its idle state.}
\end{figure}

\begin{figure}[H]
  \centering
  % \includegraphics[width=0.5\textwidth]{figures/timer-running.png}
  \caption{Timer running: the face tint animates green and the live badge
  (Stack layer) is visible.}
\end{figure}

\begin{figure}[H]
  \centering
  % \includegraphics[width=0.4\textwidth]{figures/timer-name-dialog.png}
  \caption{Ending a session asks for a name; empty names cannot be saved.}
\end{figure}

\begin{figure}[H]
  \centering
  % \includegraphics[width=0.5\textwidth]{figures/timer-logs.png}
  \caption{Today's logs with durations and the daily total.}
\end{figure}

\subsection{Authentication and states}

\begin{figure}[H]
  \centering
  % \includegraphics[width=0.5\textwidth]{figures/auth-signin.png}
  \caption{Sign-in screen (cloud mode); sign-up and password reset are
  reached from here.}
\end{figure}

\begin{figure}[H]
  \centering
  % \includegraphics[width=0.5\textwidth]{figures/error-state.png}
  \caption{Error state with retry, shown when the cloud is unreachable.}
\end{figure}

\section{Conclusion}

QuestFlow demonstrates all of the required Flutter concepts --- diverse
widgets, responsive layout, navigation, animations, state management,
theming with custom reusable widgets, and validated forms --- inside one
working product with real cloud persistence, automated tests and continuous
deployment, rather than a static UI mock-up.

\end{document}
